\documentclass[aps,preprint]{revtex4}%
\usepackage{amsfonts}
\usepackage{amsmath}
\usepackage{amssymb}
\usepackage{graphicx}%
\setcounter{MaxMatrixCols}{30}
%TCIDATA{OutputFilter=latex2.dll}
%TCIDATA{Version=4.00.0.2321}
%TCIDATA{CSTFile=revtex4.cst}
%TCIDATA{Created=Friday, December 29, 2006 11:17:19}
%TCIDATA{LastRevised=Thursday, February 01, 2007 09:19:37}
%TCIDATA{<META NAME="GraphicsSave" CONTENT="32">}
%TCIDATA{<META NAME="DocumentShell" CONTENT="Articles\SW\REVTeX 4">}
%TCIDATA{Language=American English}
\newtheorem{theorem}{Theorem}
\newtheorem{acknowledgement}[theorem]{Acknowledgement}
\newtheorem{algorithm}[theorem]{Algorithm}
\newtheorem{axiom}[theorem]{Axiom}
\newtheorem{claim}[theorem]{Claim}
\newtheorem{conclusion}[theorem]{Conclusion}
\newtheorem{condition}[theorem]{Condition}
\newtheorem{conjecture}[theorem]{Conjecture}
\newtheorem{corollary}[theorem]{Corollary}
\newtheorem{criterion}[theorem]{Criterion}
\newtheorem{definition}[theorem]{Definition}
\newtheorem{example}[theorem]{Example}
\newtheorem{exercise}[theorem]{Exercise}
\newtheorem{lemma}[theorem]{Lemma}
\newtheorem{notation}[theorem]{Notation}
\newtheorem{problem}[theorem]{Problem}
\newtheorem{proposition}[theorem]{Proposition}
\newtheorem{remark}[theorem]{Remark}
\newtheorem{solution}[theorem]{Solution}
\newtheorem{summary}[theorem]{Summary}
\newenvironment{proof}[1][Proof]{\noindent\textbf{#1.} }{\ \rule{0.5em}{0.5em}}
\begin{document}
\preprint{ }
\title{Approximate First Order Mie Coefficients}
\author{Brian Weiner}
\affiliation{Dept of Physics, Penn State Univ.}
\startpage{1}

\section{Approximate First Order Mie Coefficients}

The following is extracted from Mehboob Alam \& Yehia Massoud, "A Closed-Form
Analytical Model for Single Nanoshells", IEEE Trans. on NanoTechnology, Vol 5,
No. 3, May 2006.

The permeabilities are set as $\mu_{1}=\mu_{2}=\mu_{3}=1.$ The arguments of
the Bessel functions are%
\begin{align}
y  &  =\frac{b\omega}{c}=k_{0}b;\\
x  &  =\frac{a\omega}{c}=k_{0}a.
\end{align}
The relative indices of refraction are given by%
\begin{equation}
n_{1}^{2}=\frac{\varepsilon_{1}}{\varepsilon_{3}};n_{2}^{2}=\frac
{\varepsilon_{2}\left(  \omega\right)  }{\varepsilon_{3}},
\end{equation}
where $\varepsilon_{3}$ is the permittivity of the medium; $\varepsilon_{2}$
the permittivity of the shell and $\varepsilon_{1}$ is the permittivity of the
core. The first order (dipole) electrical Mie coefficient is
\begin{equation}
a_{1}=\frac{\frac{4y^{3}}{5}a_{1a}\left(  x\right)  a_{1ay}\left(
y,n_{2}\right)  +a_{1b}\left(  x\right)  a_{1by}\left(  y,n_{2}\right)
}{a_{1a}\left(  x\right)  a_{1cy}\left(  y,n_{2}\right)  -4a_{1b}\left(
x\right)  a_{1dy}\left(  y,n_{2}\right)  },
\end{equation}
where $a_{1a}\left(  x\right)  $ is given by
\begin{align}
a_{1a}\left(  x\right)   &  =80\left(  n_{1}^{2}+2n_{2}^{2}\right)  -8\left(
n_{1}^{2}-n_{2}^{2}\right)  \left(  n_{1}^{2}+10n_{2}^{2}\right)
x^{2}\nonumber\\
&  +2n_{2}^{2}\left(  2n_{1}^{4}+7n_{1}^{2}n_{2}^{2}-10n_{2}^{4}\right)
x^{4}-n_{1}^{2}n_{2}^{2}\left(  3n_{1}^{2}-4n_{2}^{2}\right)  x^{6}\nonumber\\
&  \simeq a_{1a0}+a_{1a2}x^{2}+a_{1a4}x^{4}%
\end{align}
$a_{1b}\left(  x\right)  $ is given by%
\begin{align}
a_{1b}\left(  x\right)   &  =\left(  n_{1}^{2}-n_{2}^{2}\right)  \left(
n_{1}^{2}n_{2}^{2}x^{7}-5\left(  n_{1}^{2}+n_{2}^{2}\right)  x^{5}%
+50x^{3}\right) \nonumber\\
&  =a_{1b3}x^{3}+a_{1b5}x^{5}+a_{1b7}x^{7}\simeq a_{1b3}x^{3}+a_{1b5}x^{5}%
\end{align}
$a_{1ay}\left(  y,n_{2}\right)  $ is given by%
\begin{align}
a_{1ay}\left(  y,n_{2}\right)   &  =\left(  n_{2}^{2}-1\right)  n_{2}^{2}%
y^{7}-\left(  n_{2}^{2}-1\right)  \left(  5\left(  n_{2}^{2}+1\right)
+50\right)  y^{3}\nonumber\\
&  =a_{1ay3}y^{3}+a_{1ay7}y^{7}\simeq a_{1ay3}y^{3}%
\end{align}
$a_{1by}\left(  y,n_{2}\right)  $ is given by%
\begin{align}
&  a_{1by}\left(  y,n_{2}\right)  =80\left(  1+n_{2}^{2}\right)  +8\left(
n_{2}^{2}-1\right)  \left(  10n_{2}^{2}+1\right)  y^{2}\nonumber\\
&  \qquad\qquad\qquad\qquad-2n_{2}^{2}\left(  10n_{2}^{4}-7n_{2}^{2}-2\right)
y^{4}+n_{2}^{4}\left(  4n_{2}^{2}-3\right)  y^{6}\nonumber\\
&  =a_{1by0}+a_{1by2}y^{2}+a_{1by4}y^{4}+a_{1by6}y^{6}\simeq a_{1by0}%
+a_{1by2}y^{2}+a_{1by4}y^{4}%
\end{align}
$a_{1cy}\left(  y,n_{2}\right)  $ is given by%
\begin{align}
&  a_{1cy}\left(  y,n_{2}\right)  =-60i\left(  n_{2}^{2}+2\right)
+\nonumber\\
&  \qquad\qquad\left(  n_{2}^{2}-1\right)  \left(  6i\left(  10+n_{2}%
^{2}\right)  y^{2}-40y^{3}\right)  -n_{2}^{2}\left(  3i\left(  n_{2}%
^{2}-4\right)  y^{4}-4\left(  n_{2}^{2}-2\right)  y^{5}\right) \nonumber\\
&  \left.  \qquad\qquad\right.  =a_{1cy0}+a_{1cy2}y^{2}+a_{1cy3}y^{3}%
+a_{1cy4}y^{4}+a_{1cy5}y^{5}%
\end{align}
and $a_{1dy}\left(  y,n_{2}\right)  $ is given by%
\begin{align}
&  a_{1dy}\left(  y,n_{2}\right)  =48i\left(  n_{2}^{2}-1\right)  +24i\left(
n_{2}^{2}-1\right)  \left(  1+n_{2}^{2}\right)  y^{2}\nonumber\\
&  +16\left(  2n_{2}^{2}+1\right)  y^{3}-6in_{2}^{4}\left(  n_{2}^{4}%
+5n_{2}^{2}-2\right)  y^{4}+8n_{2}^{2}\left(  2n_{2}^{2}-1\right)
y^{5}+3in_{2}^{4}\left(  n_{2}^{2}-3\right)  y^{6}\nonumber\\
&  \qquad\qquad\qquad\qquad\qquad\qquad\qquad\qquad\qquad\qquad\qquad
\qquad-2n_{2}^{4}\left(  2n_{2}^{2}-3\right)  y^{7}\nonumber\\
&  =a_{1dy0}+a_{1dy2}y^{2}+a_{1dy3}y^{3}+a_{1dy4}y^{4}+a_{1dy5}y^{5}%
+a_{1dy6}y^{6}+a_{1dy7}y^{7}\nonumber\\
&  \simeq a_{1dy0}+a_{1dy2}y^{2}+a_{1dy3}y^{3}+a_{1dy4}y^{4}+a_{1dy5}y^{5}.
\end{align}
Expanding the products that occur in the Mie coefficients up to order 5 gives
\begin{equation}
\frac{4y^{3}}{5}a_{1a}\left(  x\right)  a_{1ay}\left(  y,n_{2}\right)
\simeq0\qquad\qquad\qquad\qquad\qquad\qquad\qquad\qquad\qquad
\end{equation}%
\begin{align}
&  a_{1b}\left(  x\right)  a_{1by}\left(  y,n_{2}\right)  \simeq\left\{
a_{1b3}x^{3}+a_{1b5}x^{5}\right\}  \times\left\{  a_{1by0}+a_{1by2}%
y^{2}+a_{1by4}y^{4}\right\} \nonumber\\
&  \left.  \qquad\qquad\qquad\right.  \simeq\left\{  a_{1b3}x^{3}+a_{1b5}%
x^{5}\right\}  a_{1by0}+a_{1b3}a_{1by2}x^{3}y^{2}%
\end{align}%
\begin{align}
&  a_{1a}\left(  x\right)  a_{1cy}\left(  y,n_{2}\right)  =\nonumber\\
&  \left\{  a_{1a0}+a_{1a2}x^{2}+a_{1a4}x^{4}\right\}  \left\{  a_{1cy0}%
+a_{1cy2}y^{2}+a_{1cy3}y^{3}+a_{1cy4}y^{4}+a_{1cy5}y^{5}\right\} \nonumber\\
&  \simeq a_{1a0}\left\{  a_{1cy0}+a_{1cy2}y^{2}+a_{1cy3}y^{3}+a_{1cy4}%
y^{4}+a_{1cy5}y^{5}\right\} \nonumber\\
&  \qquad\qquad\qquad\qquad\qquad\qquad+a_{1a2}x^{2}\left\{  a_{1cy2}%
y^{2}+a_{1cy3}y^{3}\right\}  +a_{1a4}a_{1cy0}x^{4}%
\end{align}%
\begin{align}
&  a_{1b}\left(  x\right)  a_{1dy}\left(  y,n_{2}\right)  =\left\{
a_{1b3}x^{3}+a_{1b5}x^{5}+a_{1b7}x^{7}\right\}  \times\nonumber\\
&  \left\{  a_{1dy0}+a_{1dy2}y^{2}+a_{1dy3}y^{3}+a_{1dy4}y^{4}+a_{1dy5}%
y^{5}+a_{1dy6}y^{6}+a_{1dy7}y^{7}\right\} \nonumber\\
&  \left.  \qquad\qquad\qquad\right.  \simeq a_{1b3}x^{3}\left\{
a_{1dy0}+a_{1dy2}y^{2}\right\}  +a_{1dy0}a_{1b5}x^{5}.
\end{align}
This gives (in terms of the power series expansion coefficients) an
approximate first order electrical Mie coefficient
\begin{align}
a_{1}  &  \simeq\left\{  \left\{  a_{1b3}x^{3}+a_{1b5}x^{5}\right\}
a_{1by0}+a_{1b3}a_{1by2}x^{3}y^{2}\right\}  \left\{  a_{1a0}a_{1cy0}%
+a_{1a0}a_{1cy2}y^{2}\right\}  ^{-1}\nonumber\\
&  =\frac{1}{a_{1a0}a_{1cy0}}\left\{  \left\{  a_{1b3}x^{3}+a_{1b5}%
x^{5}\right\}  a_{1by0}+a_{1b3}a_{1by2}x^{3}y^{2}\right\}  \left\{
1-\frac{a_{1cy2}}{a_{1cy0}}y^{2}\right\} \nonumber\\
&  =\frac{1}{a_{1a0}a_{1cy0}}\left\{  \left\{  a_{1b3}x^{3}+a_{1b5}%
x^{5}\right\}  a_{1by0}+a_{1b3}a_{1by2}x^{3}y^{2}\right\}  -\frac{a_{1cy2}%
}{a_{1cy0}^{2}a_{1a0}}x^{3}y^{2}\nonumber\\
&  =\frac{a_{1by0}}{a_{1a0}a_{1cy0}}\left\{  a_{1b3}x^{3}+a_{1b5}%
x^{5}\right\}  +\left\{  \frac{a_{1b3}a_{1by2}}{a_{1a0}a_{1cy0}}%
-\frac{a_{1cy2}}{a_{1cy0}^{2}a_{1a0}}\right\}  x^{3}y^{2}.
\end{align}
Expanding this coefficient in terms of the indices of refraction we obtain the
approximate first order electric Mie coefficient
\begin{align}
&  \frac{a_{1by0}}{a_{1a0}a_{1cy0}}\left\{  a_{1b3}x^{3}+a_{1b5}x^{5}\right\}
+\left\{  \frac{a_{1b3}a_{1by2}}{a_{1a0}a_{1cy0}}-\frac{a_{1cy2}}{a_{1cy0}%
^{2}a_{1a0}}\right\}  x^{3}y^{2}\nonumber\\
&  =i\frac{\left(  1+n_{2}^{2}\right)  \left(  n_{1}^{2}-n_{2}^{2}\right)
}{\left(  n_{1}^{2}+2n_{2}^{2}\right)  60\left(  n_{2}^{2}+2\right)  }\left\{
50x^{3}-5\left(  n_{1}^{2}+n_{2}^{2}\right)  x^{5}\right\} \nonumber\\
&  -\left\{  \frac{50\left(  n_{1}^{2}-n_{2}^{2}\right)  8\left(  n_{2}%
^{2}-1\right)  \left(  10n_{2}^{2}+1\right)  }{80\left(  n_{1}^{2}+2n_{2}%
^{2}\right)  60i\left(  n_{2}^{2}+2\right)  }+\frac{\left(  n_{2}%
^{2}-1\right)  6i\left(  10+n_{2}^{2}\right)  }{\left[  -60i\left(  n_{2}%
^{2}+2\right)  \right]  ^{2}80\left(  n_{1}^{2}+2n_{2}^{2}\right)  }\right\}
x^{3}y^{2}\nonumber\\
&  =i\frac{\left(  1+n_{2}^{2}\right)  \left(  n_{1}^{2}-n_{2}^{2}\right)
}{\left(  n_{1}^{2}+2n_{2}^{2}\right)  60\left(  n_{2}^{2}+2\right)  }\left\{
50x^{3}-5\left(  n_{1}^{2}+n_{2}^{2}\right)  x^{5}\right\} \nonumber\\
&  -\left\{  \frac{\left(  n_{1}^{2}-n_{2}^{2}\right)  \left(  n_{2}%
^{2}-1\right)  \left(  10n_{2}^{2}+1\right)  }{\left(  n_{1}^{2}+2n_{2}%
^{2}\right)  12i\left(  n_{2}^{2}+2\right)  }+\frac{\left(  n_{2}%
^{2}-1\right)  \left(  10+n_{2}^{2}\right)  }{\left[  10\left(  n_{2}%
^{2}+2\right)  \right]  ^{2}80\left(  n_{1}^{2}+2n_{2}^{2}\right)  }\right\}
x^{3}y^{2}.
\end{align}
The approximate first order magnetic Mie coefficient is given as
\begin{equation}
b_{1}=\frac{i}{45}\left\{  \left(  1-n_{2}^{2}\right)  y^{5}-3n_{1}%
\frac{\left(  n_{1}^{2}-n_{2}^{2}\right)  }{\left(  1+2n_{1}\right)  }%
x^{5}\right\}  .
\end{equation}



\end{document}